\section{Benchmark Construction and Reproducibility}
\label{app:benchmarking}

% Question: What detector data form the benchmark, and how is the synthetic transit added?
% Answer: A known transit is projected into a real out-of-transit NIRISS/SOSS time series while preserving the observed background, contaminants, and noise.
% Evidence: Injector manifest and the archived benchmark release.
\subsection{Detector-level benchmark}

The benchmark is constructed from out-of-transit observations of WASP-17 with
NIRISS/SOSS. It contains 229 Stage-2 \texttt{rateints} integrations. Synthetic
transit modulation is applied to integrations 51 through 177, inclusive, so
the injected event contains 127 integrations and the remaining 102 form its
complement. The delivered science values, uncertainties, and data-quality
flags are retained, and no background subtraction is undone or reapplied.

Let $Y^{\mathrm{OOT}}_{tp}$ denote the observed out-of-transit detector time
series, $\bm{s}_{o,t}$ the denoised expectation of the target source for order
$o$, $\bm{T}^{\mathrm{true}}_{o,t}$ the known injected transit factor, and
$\bm{A}_o$ the corresponding detector projection operator. The injected cube
is
\begin{equation}
 Y^{\mathrm{inj}}_{tp}=Y^{\mathrm{OOT}}_{tp}
 +\sum_o\left[\bm{A}_o\left\{\bm{s}_{o,t}\odot
 \left(\bm{T}^{\mathrm{true}}_{o,t}-\bm{1}\right)\right\}\right]_p .
 \label{eq:v45-injector}
\end{equation}
Only the target-source expectation is modulated. The realised detector noise,
additive background, field sources, higher-order structure, and temporal
variability remain as observed. The source expectation is denoised before the
transit is applied so that the injector does not attenuate a particular noise
realisation. Contributions from the two spectral orders are constructed
separately and summed in detector space, including where their supports
overlap.

The projection operators use the PASTASOSS trace and wavelength solution
\citep{BainesEtAl2023Trace,BainesEtAl2023Wavelength} and the WebbKernel
implementation of the ATOCA response mapping
\citep{DarveauBernierEtAl2022,JWSTPipeline2025}. Their support logic fails
closed: empirical target light without a valid model response causes the
injection to stop rather than being clipped or assigned a unity transit
factor. This prevents unsupported source flux from being silently left
unmodulated.

% Question: How is information about the injected spectrum kept out of the recovery methods?
% Answer: Injection, delivery, recovery, and scoring use separate inputs, and the truth is opened only after each detector prediction is final.
% Evidence: Benchmark manifests and the recovery and scoring interfaces.
\subsection{Separation of injection, recovery, and scoring}

The experiment separates four stages. The injector owns the atmospheric
truth, its source model, and its detector projection. Delivery provides one
opaque detector cube and the same declared calibration products to every
recovery method. Each method independently constructs its calibrations,
configuration, masks, optimiser state, diagnostics, and final detector
prediction. Scoring compares the recovered and injected spectra only after
that prediction has been finalised. The injector cannot inspect a recovery
aperture, spatial-profile estimate, recovered spectrum, or truth residual.

Injected truth is therefore unavailable when constructing the detector
response, masks, nuisance basis, starting points, model-selection rules, or
convergence tests. A completed recovery records the recovered spectrum on each
declared grid, the detector prediction, convergence status, configuration and
source hashes, and a timestamped manifest. A fit that does not satisfy the
predefined numerical criteria may be retained as a diagnostic, but it does not
enter the primary comparison. Any method change made after inspecting the
truth must be evaluated on a new hidden case before it can support a performance
claim.

The same rule applies to the planned repair of the known spatial-response
defect. One global correction family must be applied to all columns, calibrated
on alternating out-of-transit folds, and selected by held-out detector
likelihood. It must preserve flux and the order-specific and aggregate
acceptance criteria, and it must be fixed before the recovered spectrum is
compared with the truth. The known defect can motivate the model family to be
tested, but its measured recovery error cannot select the correction.

% Question: Which wavelengths and metrics enter the common comparison?
% Answer: Every method is scored on complete R=100 bins within the direct PASTASOSS calibration domain.
% Evidence: Calibrated-domain mask and common-grid manifest.
\subsection{Common reporting support}

All methods are compared on the same physical $R=100$ grid and with the same
predefined support mask. The order-1 PASTASOSS training anchors for this visit
end at detector coordinate $x=73.62893$, corresponding to
$2.75977449\,\mu\mathrm{m}$. Detector columns 38 through 73 depend on
polynomial extrapolation and are therefore excluded from the primary
comparison. Any reporting bin with even partial dependence on these columns is
removed in full. This leaves 141 complete bins, with the final retained bin
ending at approximately $2.740018\,\mu\mathrm{m}$. The mask is determined by
the calibration domain, applied identically to all methods, and used only for
reporting rather than fitting.

The primary accuracy statistic is the unweighted absolute root-mean-square
error on these common bins. Mean bias and morphology error are reported
separately so that an achromatic offset is not confused with an incorrect
spectral shape. Order-specific, overlap, crossing-region, and supported-red
diagnostics accompany the aggregate score. Weighted error is treated as a
sensitivity statistic. Zero-injection, flat-depth, and noiseless engineering
tests are reported separately from the main recovery experiment.

% Question: Can every comparison pipeline reproduce its normal preprocessing from an integration-level cube?
% Answer: No. A matched group-stage experiment is required to isolate the effect of preprocessing performed before ramp fitting.
% Evidence: Injector fairness programme and current comparison status.
\subsection{Integration-level limitation and paired group-stage experiment}

The present benchmark begins with integration-level count-rate images. It does
not contain the nondestructive detector groups from which those images were
fitted. Pipelines such as Ahsoka and supreme-SPOON/exoTEDRF normally perform
characteristic $1/f$ corrections before or during ramp fitting and cannot
reproduce that step from the delivered cube. The current benchmark can compare
recoveries from identical integration-level data, but it cannot determine how
group-stage preprocessing changes their relative performance.

A paired sensitivity experiment is therefore planned. One common group-level
baseline will be used to create matched group-stage and integration-stage
injections from the same physical source deficit. Each pipeline will retain
its own predefined preprocessing, and NOVA will also be run after the same
group-stage $1/f$ correction. The comparison will include zero-injection
identity, the delivered detector deficit, white-light scatter, red-end
recovery, and common-bin spectral error. This experiment is prospective and
does not form part of the present result.

A preliminary scene-aware injector is also excluded from the current
benchmark. Its broad full-column masks remove a material fraction of the
target signal together with contaminants. Before such a model can be used, it
must be replaced by out-of-transit scene modelling that protects the target
trace and quantifies any remaining contaminant leakage.

% Question: What additional experiments are needed before the comparison can be generalised?
% Answer: Performance must remain stable across detector-response mismatch, visits, atmospheric spectra, and noise realisations.
% Evidence: Benchmark limitations and planned validation programme.
\subsection{Structural neutrality and ensemble tests}

The injector and NOVA use related families of trace geometry, wavelength
mapping, and spectral mixing. They do not share fitted detector responses or
the same executable implementation, but their common modelling ancestry creates
an inverse-crime risk. The benchmark programme will therefore include hidden
injections with deliberately mismatched spatial profiles, wavelength maps,
line-spread functions, throughput, background morphology, and temporal
systematics.

WASP-17 is a development visit rather than sufficient validation by itself.
Additional visits must pass a fixed admission audit covering target and visit
identity, observing mode, out-of-transit coverage, integration chronology,
data quality, background structure, field sources, and order-3 morphology. The
same injected spectrum will first be recovered across the admitted visits,
followed by additional spectral morphologies and a method-neutral noise
ensemble. At least one visit will remain held out from model-branch selection.
Until these tests and the comparison-pipeline matrix are complete, the current
benchmark supports a controlled single-visit result rather than a universal
ranking of reduction methods.

% Question: What record is required to reproduce or audit a benchmark result?
% Answer: The data, code, configuration, prediction, and scoring identities are preserved for every reported recovery.
% Evidence: Release manifest and portable handoff audit.
\subsection{Provenance and minimum report}

For byte-level provenance, the archived integration-level cube and reporting
mask are identified internally as V45 and R7, respectively. These identifiers
describe repository lineage rather than scientific components of the method.
Each release records content hashes for the code, configuration, calibration
and response products, source and background products, detector arrays, masks,
final predictions, and score tables. Large cluster-only arrays may be
represented in the portable archive by authenticated manifests. Superseded,
diagnostic, unsuccessful, and prospective branches retain an explicit status
so that they cannot be mistaken for the reported method.

\begin{table}[t]
 \centering
 \footnotesize
 \caption{Minimum record for each benchmark recovery.}
 \label{tab:benchmark-report}
 \setlength{\tabcolsep}{3pt}
 \begin{tabularx}{0.97\columnwidth}{@{}>{\raggedright\arraybackslash}p{0.29\columnwidth}>{\raggedright\arraybackslash}X@{}}
  \toprule
  Category & Required record \\
  \midrule
  Delivery & Cube, uncertainties, flags, time grid, wavelength support, and
  response and configuration identities \\
  Numerical validity & Termination state, final-weight confirmation,
  multistart agreement, rank, conditioning, active bounds, and physical checks \\
  Spectral accuracy & Absolute and morphology error, mean bias, regional
  residuals, and separate order products where available \\
  Uncertainty & Standardised residuals and spectral covariance for each
  recovery; empirical coverage across the planned noise ensemble \\
  Detector diagnostics & Residual and robust-weight summaries by time, order,
  detector region, and nuisance component \\
  Scope & Primary, diagnostic, unsuccessful, or prospective status, together
  with the claims that the case can and cannot support \\
  \bottomrule
 \end{tabularx}
\end{table}
